\section{Finite discs with congruence fibers}\label{sec:geometry}

We construct finite graphs $X_j$ containing distinguished boundary cycles
$C_j$. The goal of this section is a uniform statement about short loops:
at any fixed length bound, loops sufficiently far out in the family have
zero signed crossing count for a fixed ray from the center of the underlying
disc. This count agrees with winding whenever the projected loop avoids
the center.

\subsection{Congruence quotients and a separated holonomy}
We need a quotient tower with uniformly bounded successive kernels, so that
lifting base edges preserves bounded degree. We will also choose compatible
twists whose trace separation rules out nonzero crossing count on short
closed paths in sufficiently late graphs $X_j$. The choice of $\SL_3(\Z)$ supplies
both this quotient tower and the uniform averaging estimates used in
Section~\ref{sec:transfer}.

Let $\Gamma=\SL_3(\Z)$ and let $S$ be the finite symmetric generating set of
elementary matrices $I\pm E_{ab}$, $a\ne b$. To see that these generate,
apply the Euclidean algorithm by determinant-one row operations to the
primitive first column of an integral determinant-one matrix, then to its
remaining $2$-by-$2$ block. Each integral elementary operation is a power
of a member of $S$. Put
\[
 Q_0=\{1\},\qquad Q_k=\SL_3(\Z/2^k\Z)\quad(k\ge1),
\]
with their reduction maps. All reductions from $\Gamma$, and between
successive $Q_k$, are surjective. Indeed, over $\Z/2^k\Z$ an invertible
matrix has a unit in its first column. Elementary row operations move it to
the first position and clear the remaining entries. Continuing reduces the
matrix to a diagonal one. The identity
\[
 \begin{pmatrix}0&u\\-u^{-1}&0\end{pmatrix}
 =\begin{pmatrix}1&u\\0&1\end{pmatrix}
  \begin{pmatrix}1&0\\-u^{-1}&1\end{pmatrix}
  \begin{pmatrix}1&u\\0&1\end{pmatrix}
\]
expresses the needed determinant-one row interchanges and diagonal
matrices by elementary matrices. Every elementary matrix over the quotient
lifts to an integral one, itself generated by $S$.
There is a uniform bound
\begin{equation}\label{eq:kernelbound}
 h:=2^9\ \ge\ \abs{\ker(Q_{k+1}\to Q_k)}\quad(k\ge0).
\end{equation}
For $k\ge1$, a kernel element has the form $I+2^kA$ modulo $2^{k+1}$ and is
specified by the nine binary entries of $A$; the same bound holds for $Q_1$.

Choose a $2$-adic unit $t\in1+2\Z_2$ transcendental over $\Q$. Such a choice
exists because $1+2\Z_2$ is uncountable and the algebraic elements over
$\Q$ are countable. For every nonzero integer $m$,
\begin{equation}\label{eq:transcendentaltrace}
 t^m+t^{-m}+1\notin\Z.
\end{equation}
Otherwise multiplying the alleged equality by $t^{\abs m}$ gives a nonzero
polynomial over $\Z$ satisfied by $t$.
Define
\[
 g_k=\operatorname{diag}(t,t^{-1},1)\pmod{2^k}\in Q_k\quad(k\ge1),
 \qquad g_0=1.
\]
These elements are compatible with reduction.

\begin{lemma}[Trace separation]\label{lem:trace}
For each integer $A\ge1$ there is an integer $N(A)\ge1$ such that, whenever
$\ell\ge N(A)$, $0<\abs m\le A$, and $w\in\Gamma$ has word length at most
$A$ with respect to $S$,
\[
 \tr(g_\ell^m)\ne\tr(w)\pmod{2^\ell}.
\]
\end{lemma}
\begin{proof}
There are only finitely many pairs $(m,w)$ in question. Each difference
$t^m+t^{-m}+1-\tr(w)$ is a nonzero $2$-adic integer by
\eqref{eq:transcendentaltrace}. Choose $N(A)$ greater than all their
$2$-adic valuations.
\end{proof}

\subsection{The base discs and the lifted edges}
Use the triangular lattice with coordinates $(a,b)\in\Z^2$ and neighbor
differences
\[
 \pm(1,0),\quad\pm(0,1),\quad\pm(1,-1).
\]
Embed its two coordinate axes at angle $\pi/3$. Set
\[
 \rho(a,b)=\max\{\abs a,\abs b,\abs{a+b}\}.
\]
For $j\ge2$, let $D_j$ be the finite induced graph on $\rho\le j$,
equipped with its graph metric. Its elementary triangles fill a closed
topological disc. The boundary $C_j$, given by $\rho=j$, is a cycle of
length $6j$; we equip its vertex set with its own cycle metric. When used as a
topological circle, $C_j$ denotes the geometric realization of this cycle.
The depth of a vertex
$u\in D_j$ is
\[
 k_j(u)=j-\rho(u).
\]
Depths of neighboring vertices differ by at most one.

Fix a ray from the origin meeting no other lattice vertex. For an oriented
lattice edge $(u,v)$ not incident to the origin, let $a(u,v)\in\{-1,0,1\}$
be its signed intersection with this ray. Set $a(u,v)=0$ for edges incident
to the origin. Choose the sign so that the positively oriented boundary
has total charge one. Then $a(v,u)=-a(u,v)$, and the sum of $a$ on any
closed edge path avoiding the origin equals its winding number.

The vertices of $X_j$ are
\[
 \{(u,q):u\in D_j,\ q\in Q_{k_j(u)}\}.
\]
Put undirected edges of length one as follows; repeated edges and loops
may be discarded.
\begin{enumerate}
\item Within the fiber over $u$, join $q$ to $sq$ for every $s\in S$.
\item If $(u,v)$ is an oriented base edge with equal depth $k$, join
$(u,q)$ to $(v,qg_k^{a(u,v)})$.
\item If $(u,v)$ goes from depth $k+1$ to depth $k$, join
$(u,q)$ to $(v,\bar qg_k^{a(u,v)})$, where $\bar q$ denotes reduction.
The reverse direction uses the inverse edge relation.
\end{enumerate}
Left Cayley multiplication and right holonomy commute. This choice of
sides will allow a short closed path to yield a conjugacy relation in a
single quotient.

Each $X_j$ is connected: the fibers are connected Cayley graphs, and every
base edge has nonempty relations to both endpoint fibers. Each vertex has
degree at most
\begin{equation}\label{eq:degree}
 \Delta=\abs S+6h.
\end{equation}
Indeed, a base neighbor gives a bijection, a single parent, or at most $h$
children. Boundary fibers are trivial, so the cycle $C_j$ includes in $X_j$
by $u\mapsto(u,1)$. Figure~\ref{fig:disc} shows the depth labels and
the two types of inter-fiber edge relation.

These graphs do have uniformly short cycles. In a fiber of depth $k\ge1$,
the commuting generators $A=I+E_{12}$ and $B=I+E_{13}$ give the four
distinct vertices $q,Aq,BAq,Bq$ of a Cayley square. Thus the short-loop
statement below concerns their charge, not the absence of short cycles.

\begin{figure}[t]
\centering
\begin{tikzpicture}[scale=.77,>=Stealth,font=\small]
 \foreach \r in {1,2,3} {
  \draw[gray!65] (0:\r) -- (60:\r) -- (120:\r) -- (180:\r) -- (240:\r) -- (300:\r) -- cycle;
 }
 \foreach \ang in {0,60,...,300} {\draw[gray!35] (0,0)--(\ang:3);}
 \draw[very thick,blue!55!black] (0:3)--(60:3)--(120:3)--(180:3)--(240:3)--(300:3)--cycle;
 \draw[dashed,thick,->] (0,0)--(19:3.65) node[right]{cut};
 \fill (0,0) circle (2pt) node[below left] {$Q_3$};
 \foreach \r/\q in {1/2,2/1,3/0} {
  \fill (120:\r) circle (2pt);
  \node[above left] at (120:\r) {$Q_{\q}$};
 }
 \node[blue!55!black] at (0,-3.05) {$C_j$: singleton fibers};
 \node at (5.9,2.85) {equal depth $k$};
 \foreach \y in {1.5,1.95,2.4} {
  \fill (4.7,\y) circle (2pt);
  \fill (7.1,\y) circle (2pt);
  \draw[->,gray!75] (4.8,\y)--(7,\y);
 }
 \node at (5.9,1) {$q\longmapsto qg_k^a$};
 \node at (4.7,-.25) {$Q_k$};
 \node at (7.1,-.25) {$Q_{k+1}$};
 \fill (4.7,-1.2) circle (2pt);
 \foreach \y in {-.7,-1.2,-1.7} {
  \fill (7.1,\y) circle (2pt);
  \draw[gray!75] (4.8,-1.2)--(7,\y);
 }
 \node at (5.9,-2.1) {$\bar z=qg_k^a$};
 \node at (5.9,-2.8) {within a fiber: $q\longleftrightarrow sq$};
\end{tikzpicture}

\caption{The base disc and its fibers, shown schematically for $j=3$.
The shell at depth $k$ carries $Q_k$; the boundary carries $Q_0=\{1\}$.
Crossing the dashed ray applies right multiplication by $g_k^{a(u,v)}$,
with reduction to the shallower level when the depths differ. The center
has depth $j$. The insets show selected inter-fiber edges: a twisted
matching at equal depth and a shallow vertex with its deep children.
In each inset, $a=a(u,v)$ for the base edge oriented from the left-hand
fiber to the right-hand fiber.
The number of drawn vertices is schematic, not the quotient cardinality.}
\label{fig:disc}
\end{figure}

\subsection{Separated unions and their boundary maps}
Choose $b_j\in C_j$. In each of the three unions
\[
 B=\coprod_{j\ge2}C_j,\qquad D=\coprod_{j\ge2}D_j,\qquad
 X=\coprod_{j\ge2}X_j,
\]
join the chosen basepoints of successive components by an auxiliary edge
of length $M_j=100+j^2$; in $X_j$ use $(b_j,1)$.
Use the induced path distance on the original vertex sets, without adding
vertices along these auxiliary edges. Inside each component this is its
original graph or cycle metric: leaving a component cannot provide a
shortcut back into it. The distances between distinct components tend to
infinity outside a finite family, so these are coarse disjoint unions.

The spaces are uniformly discrete and proper, and have bounded geometry.
To check the last assertion directly, the weighted graph has degree at
most $\Delta+2$ and every edge has length at least one. Hence a ball of
radius $R$ has at most
$\sum_{m=0}^{\lfloor R\rfloor}(\Delta+2)^m$ vertices. The same argument
applies to $B$ and $D$ with their smaller degree bounds. There are proper
$1$-Lipschitz inclusions
\begin{equation}\label{eq:boundarymaps}
 e:B\longrightarrow D,\quad u\longmapsto u,
 \qquad i:B\longrightarrow X,\quad u\longmapsto(u,1).
\end{equation}
Every path in $B$ maps to a path of the same length, including the auxiliary
links. Properness follows from injectivity and finiteness of bounded subsets
in the target. The use of the intrinsic cycle metric on $B$ ensures these
Lipschitz assertions.

For later use, at every fixed distance bound $R$ only finitely many pairs
$(u,v)\in D\times D$ from different components satisfy $d_D(u,v)\le R$.
Such a path must cross a link of length at most $R$. There are only finitely
many such links, and both endpoints of the pair lie within distance $R$
of their endpoints. Those balls are finite.

\subsection{Short loops have zero charge far out}
For an edge path in $X_j$, define its charge to be the sum of $a(u,v)$ on
its projected base steps; fiber steps contribute zero. Charges are additive
under concatenation and change sign under reversal.

\begin{lemma}[Small-loop lemma]\label{lem:smallloops}
For each integer $A\ge1$ there is $J(A)$ such that, for all $j\ge J(A)$,
every closed edge path in $X_j$ of length at most $A$ has charge zero.
One may take $J(A)=\max\{2,N(A)+\lfloor A/2\rfloor\}$, with $N(A)$ from
Lemma~\ref{lem:trace}.
\end{lemma}
\begin{proof}
Consider a closed path of length $n\le A$ and projected vertices in $D_j$.
Suppose its charge $m$ is nonzero. We first show that every projected vertex
$u$ satisfies $\rho(u)\le\lfloor n/2\rfloor$.
If not, choose a linear functional
$f\in\{\pm a,\pm b,\pm(a+b)\}$ attaining $\rho(u)$.
Every other vertex is reachable from $u$ along a shorter arc of the closed
path of length at most $\lfloor n/2\rfloor$. Since $f$ changes by at most
one in each step, including the stationary projections of fiber steps,
the entire polygon lies in the open half-plane $f>0$. This argument applies
to each occurrence of a vertex, so repeated vertices and edges cause no
change. The polygon is contractible in that half-plane, avoids the origin,
and has winding zero, contradicting its charge. In particular, the same
argument covers paths whose projection initially might have met the origin:
the supposed vertex outside the bound would force the whole projection to
avoid it.

Let $\ell$ be the minimum depth along the path. The preceding bound gives
$\ell\ge j-\lfloor A/2\rfloor$. Reduce every fiber coordinate along the
path to $Q_\ell$. A fiber step becomes left multiplication by an element of
$S$. An oriented base step becomes right multiplication by
$g_\ell^{a(u,v)}$. This is immediate in the equal-depth and deep-to-shallow
cases. In the reverse case the defining edge relation says
$q_{\mathrm{shallow}}=\bar q_{\mathrm{deep}}g_k^{-a}$, so that
$\bar q_{\mathrm{deep}}=q_{\mathrm{shallow}}g_k^a$, as required.

Commuting left and right multiplications and using closure of the path,
we obtain
\[
 w z g_\ell^m=z
\]
for its starting coordinate $z\in Q_\ell$ and a word $w$ of length at most
$A$ in $S$. Consequently $g_\ell^m=z^{-1}w^{-1}z$ and
$\tr(g_\ell^m)=\tr(w^{-1})$ modulo $2^\ell$.
Here $0<\abs m\le A$. If $j\ge J(A)$, then $\ell\ge N(A)$ and this
contradicts Lemma~\ref{lem:trace} applied to $w^{-1}$.
\end{proof}

The qualification ``$j$ sufficiently large'' is part of the conclusion.
It is precisely what is needed for a test at each fixed Rips scale; no
single short-loop bound is asserted simultaneously for all components.
